\subsection{局部凸空间族的拓扑直和}

定义 2. --- 设 E 是局部凸空间族 \((\mathrm{E}_{\iota})_{\iota \in \mathbf{I}}\) 的作为向量空间的直和 (A, II, §1.6)。对每个 \(\iota \in \mathbf{I}\) ，设 \(f_{\iota}\) 是 \(E_{t}\) 到 E 中的典范单射。我们把族 ( \(E_{t}\) ) 的拓扑直和理解为带有使每个 \(f_{t}\) 都连续的最细局部凸拓扑的空间 E（这个拓扑称为诸 \(E_{t}\) 的拓扑的直和）。

在本节余下部分，我们沿用 def. 2 的记号（除非另有明确声明），并把每个 \(E_{t}\) 借助 \(f_{t}\) 典范地等同于 E 的一个子空间。

由 II, 第 28 页给出的关于终局部凸拓扑的邻域的一般描述，我们可以按下述方式得到 E 中直和拓扑的 0 的一个基本邻域系；对每个族 \((\mathrm{V}_{\iota})_{\iota \in I}\) （其中 \(\mathrm{V}_{\iota}\) 是 \(\mathrm{E}_{\iota}\) 中 0 的一个对称凸邻域），考虑凸包 \(\Gamma((\mathrm{V}_{\iota}))\) ，即诸 \(\mathrm{V}_{\iota}\) 的并；诸 \(\Gamma((\mathrm{V}_{\iota}))\) （对所有族 \((\mathrm{V}_{\iota})\) ，或只取 \(\mathrm{V}_{\iota}\) （每个 \(\iota\) 的）属于 \(\mathrm{E}_{\iota}\) 中 0 的一个基本邻域系）构成 E 中 0 的一个基本邻域系。

例. --- 设 \((a_{\mathfrak{t}})_{\mathfrak{t}\in\mathbf{I}}\) 是向量空间 E 的一个基，并在每条直线 \(Ra_{\mathfrak{t}}\) 上考虑典范拓扑 (I, 第 2 页, 例 5)；这些拓扑的直和是 E 上最细的局部凸拓扑 (II, 第 26 页)；事实上，若 V 是 E 中的一个吸收、对称、凸集，则 \(V_{\mathfrak{t}} = V \cap Ra_{\mathfrak{t}}\) 是 \(Ra_{\mathfrak{t}}\) 中 0 的一个邻域，而 V 显然包含凸包 \(\Gamma((V_{\mathfrak{t}}))\) 。

命题 6. --- E 上的局部凸拓扑 T 是诸 \(E_{t}\) 的拓扑的直和，必须且只需以下性质成立：E 到局部凸空间 G 中的线性映射（相应地，E 上的半范数 p）连续，必须且只需对每个 \(t \in I\) ，映射 \(g \circ f_{t}\) （相应地，\(p \circ f_{t}\) ）在 \(E_{t}\) 中连续。

这是 II, 第 27 页, prop. 5 的一个特殊情形。

回忆向量空间族的直和的定义 (A, II, 第 12 页, prop. 6)，我们可以说，拓扑 \(\mathcal{T}\) 是使典范映射 \(g \mapsto (g \circ f_{\mathrm{t}})\) 为双射的唯一拓扑

\[
\mathcal {L} (\mathrm{E}; \mathrm{G}) \rightarrow \prod_ {\iota \in \mathrm{I}} \mathcal {L} (\mathrm{E} _ {\iota}; \mathrm{G})\tag{1}
\]

对每个局部凸空间 \(\mathbf{G}\) 都成立。

推论. --- 沿用 II, 第 27 页, prop. 5 的记号，假设 E 是诸 \(g_{\alpha}(\mathbf{F}_{\alpha})\) 之和。设 F 是族 \((\mathbf{F}_{\alpha})_{\alpha \in \mathbf{A}}\) 的拓扑直和，\(j_{\alpha}: \mathbf{F}_{\alpha} \to \mathbf{F}\) 是典范单射；设 \(g: \mathbf{F} \to \mathbf{E}\) 是满足 \(g \circ j_{\alpha} = g_{\alpha}\) （对所有 \(\alpha \in \mathbf{A}\) ）的线性映射。若 N 是 \(g\) 的核，则典范双射 \(\mathrm{F/N} \to \mathrm{E}\) （与 \(g\) 相伴的）是 \(\mathrm{F/N}\) 到带有拓扑 \(\mathcal{T}\) 的 E 上的一个拓扑同构。

这是 II, 第 28 页, cor. 2 的一个特殊情形，并要记住 II, 第 29 页, 例 I。

命题 7. --- 典范单射 \(j: \mathrm{E} \to \prod_{\iota \in \mathrm{I}} \mathrm{E}_{\iota}\) 在 \(\mathrm{E}\) 带有诸 \(\mathrm{E}_{\iota}\) 的直和拓扑而 \(\prod_{\iota \in \mathrm{I}} \mathrm{E}_{\iota}\) 带有乘积拓扑时连续。当 \(\mathrm{I}\) 有限时，这个映射是拓扑向量空间间的一个同构。

第一个断言来自如下事实：典范单射 \(E_{\kappa} \to \prod_{i \in I} E_{i}\) 对每个 \(\kappa \in I\) 都连续。若 I 有限，则 j 是恒等映射，于是只须证明乘积拓扑 \(T'\) 比直和拓扑 T 细。设 V 是 T 下 0 的一个凸邻域；每个集 \(V \cap E_{i}\) 都是 \(E_{i}\) 中 0 的一个凸邻域；若 n 是 I 的元素个数，则集 V 包含集 \(\frac{1}{n} \sum_{n} (V \cap E_{i})\) ，后者是 \(T'\) 下 0 的一个邻域，命题得证。

当 I 无限时，若对 I 的每个有限子集 J，用 \(E_{J}\) 表示带乘积拓扑的空间 \(\prod_{i\in J}E_{i}\) ，则 E 是诸 \(E_{J}\) （视为 E 的子空间）的归纳极限。

命题 8. --- 设 \(\mathbf{N}_{\iota}\) 是 \(\mathrm{E}_{\iota}\) 的一个子空间（对每个 \(\iota \in \mathbf{I}\) ），

\begin{enumerate}
\def\labelenumi{(\roman{enumi})}
\item
  在 \(\mathbf{N} = \sum_{\mathfrak{i}}\mathbf{N}_{\mathfrak{i}}\) 上由直和拓扑 \(\mathcal{T}\) （\(\mathbf{E}\) 上的）诱导的拓扑与诸 \(\mathbf{N}_{\mathfrak{i}}\) 的拓扑的直和相同。
\item
  典范映射 \(h\) （它把诸 \(\mathrm{E}_{\mathrm{i}} / \mathrm{N}_{\mathrm{i}}\) 的拓扑直和空间映到 \(\mathrm{E} / \mathrm{N}\) 上）(A, II, § 1.6, 公式 (26)) 是拓扑向量空间之间的一个同构。
\item
  沿用前面引进的记号，考虑 \(x = \sum_{i} \lambda_i x_i\) ，它属于 \(N \cap \Gamma((V_{t}))\) （ \((\lambda_{t})\) 是一个至多有限多个非零的数 \(\geqslant 0\) 组成的族，满足 \(\sum_{i} \lambda_{i} = 1\) ，且 \(x_{i} \in V_{i}\) （对所有 \(i \in I\) ））。
\end{enumerate}

由于诸 \(N_{i}\) 之和是直和，有 \(\lambda_{i}x_{i} \in N_{i}\) （对所有 \(i \in I\) ）；因此对所有使 \(\lambda_{i} > 0\) 的 i 也有 \(x_{i} \in N_{i} \cap V_{i}\) ，而 x 属于凸包 \(\Gamma((N_{i} \cap V_{i}))\) ，于是 (i) 得证。(ii) 把诸典范映射记作：\(f_{i}: E_{i} \to E\) ， \(h_{i}: E_{i}/N_{i} \to E/N\) ， \(p_{i}: E_{i} \to E_{i}/N_{i}\) 及 \(p: E \to E/N\) 。对每个 \(i \in I\) ，有 \(h_{i} \circ p_{i} = p \circ f_{i}\) ，命题由 II, 第 28 页, cor. 2 及第 29 页例 I 得出。

推论 1. --- 若 \(\mathbf{N}_{\iota}\) 在 \(\mathbf{E}_{\iota}\) 中是闭的（对每个 \(\iota \in \mathbf{I}\) ），则 \(\mathbf{N} = \sum \mathbf{N}_{\iota}\) 在 \(\mathbf{E}\) 中是闭的。

因为典范映射 \(p_{\iota} : \mathrm{E} \to \mathrm{E}_{\iota}\) 连续 (II, § 4.5, prop. 6)（对每个 \(\iota \in \mathbf{I}\) ），故 \(p_{\iota}^{-1}(\mathbf{N}_{\iota})\) 在 \(\mathbf{E}\) 中是闭的，从而对交 \(\mathbf{N} = \bigcap_{\iota \in \mathbf{I}} p_{\iota}^{-1}(\mathbf{N}_{\iota})\) 而言也是如此。

推论 2. --- 若每个 \(E_{t}\) 都是 Hausdorff 的，则 E 也是 Hausdorff 的，并且每个 \(E_{t}\) 在 E 中是闭的。

为证第一个断言，取 \(N_{t} = \{0\}\) （对所有 \(t \in I\) ）应用 cor. 1；为证第二个断言，取 \(N_{t} = E_{t}\) 并取 \(N_{\kappa} = \{0\}\) （对每个 \(\kappa \neq t\) ）应用 cor. 1。

我们将在 III, 第 21 页, cor. 2 中证明：若诸 \(E_{t}\) 都是 Hausdorff 的且完备的，则它们的拓扑直和 E 也是如此。

